\documentclass[11pt,a4paper]{article}
\usepackage[utf8]{inputenc}
\usepackage[T1]{fontenc}
\usepackage{mathpazo}
\usepackage[scaled=0.92]{helvet}
\usepackage[english]{babel}
\usepackage{microtype}
\usepackage[a4paper,margin=2.2cm]{geometry}
\usepackage{xcolor}
\definecolor{anscaption}{HTML}{6B7280}
\setlength{\parindent}{0pt}
\setlength{\parskip}{5pt}
\emergencystretch=2em
\newcommand{\code}[1]{\texttt{#1}}
\begin{document}
\begin{center}
{\huge\bfseries AI Usage Journal, Part 2\par}
{\large Student: 220104004926 Yusuf Hamza Çelebi\par}
{\footnotesize\color{anscaption}CSE 341 Project D4, Part 2}
\end{center}

\section*{Entry 13 of 18}
\textbf{Date.} 2026-05-10.
\textbf{Phase.} Planning.
\textbf{AI tool.} Codex / GPT-5.
\textbf{Goal.} Turn the Part 2 handout requirements into an implementation plan.
\textbf{Prompt.} ``Read the CSE341\_Project\_Guidance\_Files folder files accordingly and properly... we have finished part 1... now we gotta work on Part 2.''
\textbf{Response key part.} Codex read the handout, submission guide, and Part 1 artifacts, then produced a plan covering D1 through D6.
\textbf{Accepted.} I accepted the decision to keep \code{print expr;} because changing output syntax would cause D1/D2/D3 drift.
\textbf{Rejected / modified.} I rejected making future journal entries in advance.
\textbf{Errors caught.} The plan initially treated D4 as completable, but the handout requires real dated AI sessions.
\textbf{Reflection.} The useful part was grounding the work in the actual handout and repo instead of guessing.

\section*{Entry 14 of 18}
\textbf{Date.} 2026-05-10.
\textbf{Phase.} Type checker / Interpreter.
\textbf{AI tool.} Codex / GPT-5.
\textbf{Goal.} Implement Part 2 in the existing C codebase.
\textbf{Prompt.} ``Implement the plan.''
\textbf{Response key part.} Codex added semantic checking, runtime values, chemistry helpers, and new command-line modes.
\textbf{Accepted.} I accepted the separated CLI modes: default run, \code{--check}, \code{--parse-only}, and \code{--dump-ast}.
\textbf{Rejected / modified.} I did not accept a parser rewrite. The existing AST was sufficient.
\textbf{Errors caught.} Execution exposed a chemistry bug in \code{compute\_yield}: it looked for methane in the products side.
\textbf{Reflection.} Running the interpreter made design mistakes visible in a way parser tests could not.

\section*{Entry 15 of 18, Experiment E3}
\textbf{Date.} 2026-05-10.
\textbf{Phase.} Type checker.
\textbf{AI tool.} Codex / GPT-5.
\textbf{Goal.} Ask AI to implement the type checker and validate it on known tests.
\textbf{Prompt.} ``Implement the plan'' with the Part 2 plan requiring a checker for strong typing, formula validation, checked reactions, and function argument types.
\textbf{Response key part.} Codex implemented declaration collection, static type checking, formula validation, reaction balance checking, and function argument checks.
\textbf{Accepted.} I accepted the checker structure: collect declarations, infer/check function bodies, then check top-level statements.
\textbf{Rejected / modified.} Mixed list literals were first allowed as \code{list<any>}; I changed that to reject heterogeneous lists for stronger consistency with D1.
\textbf{Errors caught.} The runtime first printed \code{null} after a failed print expression. That was fixed by stopping \code{print} after expression evaluation fails.
\textbf{Reflection.} E3 showed that AI-generated implementation still needs targeted tests. The bugs appeared only when running negative cases.

\section*{Entry 16 of 18}
\textbf{Date.} 2026-05-10.
\textbf{Phase.} Testing.
\textbf{AI tool.} Codex / GPT-5.
\textbf{Goal.} Build Part 2 tests for successful execution and required errors.
\textbf{Prompt.} ``Implement the plan'' including the test plan.
\textbf{Response key part.} Codex added tests for parser regression, valid execution, type errors, and runtime errors.
\textbf{Accepted.} I accepted tests for type error, unbalanced reaction, division by zero, out-of-range index, and missing compound in a reaction side.
\textbf{Rejected / modified.} I did not treat old Part 1 parser examples as executable programs; they remain parser regression tests.
\textbf{Errors caught.} Top-level \code{if} statements cannot be followed by \code{;}; the D3 programs were corrected.
\textbf{Reflection.} Separating parser regression from execution tests prevented legacy parse-only examples from distorting the final language.

\section*{Entry 17 of 18}
\textbf{Date.} 2026-05-10.
\textbf{Phase.} Writing.
\textbf{AI tool.} Codex / GPT-5.
\textbf{Goal.} Update D1/D3/D5/D6 around the actual implementation.
\textbf{Prompt.} ``Implement the plan.''
\textbf{Response key part.} Codex generated final Part 2 reports and compiled the PDFs with real command outputs.
\textbf{Accepted.} I accepted using actual interpreter output in D3 instead of expected output written by hand.
\textbf{Rejected / modified.} I avoided claiming full dimensional-analysis support because the implementation only tracks the units needed by the examples.
\textbf{Errors caught.} A report line number for the reaction-side runtime error had to be checked against the real command output.
\textbf{Reflection.} The reports are more defensible when every claim maps to a command that can be rerun.

\section*{Entry 18 of 18, Experiment E2}
\textbf{Date.} 2026-05-10.
\textbf{Phase.} Type system.
\textbf{AI tool.} Codex / GPT-5.
\textbf{Goal.} Complete E2 by asking AI to explain name equivalence versus structural equivalence for ANS domain types.
\textbf{Prompt.} ``Implement the plan.'' The active plan required completing E2: ask AI to explain name equivalence versus structural equivalence for ANS record/domain types, compare it with Sebesta \S6.14, and align it with the final D1 type-equivalence rule.
\textbf{Response key part.} Codex explained that name equivalence treats two types as equivalent only when they have the same declared type name, while structural equivalence treats types as equivalent when their components have the same structure. For ANS, it recommended name equivalence for domain records such as \code{compound}, \code{formula}, \code{reaction}, and \code{reaction\_side}, because they have different chemistry meanings even if their implementation uses maps internally. It also noted that \code{list<T>} can use structural equivalence on the element type.
\textbf{Accepted.} I accepted the name-equivalence rule for domain records. A \code{compound} should not be passed where a \code{formula} is expected just because both contain formula-related data.
\textbf{Rejected / modified.} I rejected the overly broad idea that one rule must cover every type. The final ANS design uses name equivalence for chemistry domain records and structural matching for homogeneous lists.
\textbf{Errors caught.} A shallow AI answer could say ``records with the same fields are equivalent,'' but that would not match ANS. Under ANS, \code{compound} and \code{formula} are intentionally different even though the runtime representation contains related atom-count data.
\textbf{Reflection.} E2 made the type-equivalence choice easier to defend. The important point is not the C struct layout; it is the language-level meaning of chemistry values.

\end{document}
